\documentclass[12pt]{amsart}
\pagestyle{empty}
\usepackage[hmargin=1in,vmargin=1in]{geometry}
\usepackage{url,  mathrsfs}
\usepackage{hyperref}
\usepackage{fancyhdr} 
\usepackage[T1]{fontenc}


\usepackage{url,  mathrsfs}
\usepackage{hyperref}
\usepackage{amsmath}
\usepackage{graphicx,  epstopdf}
\usepackage{subfig}
\usepackage{color}
\usepackage{bbm}
\usepackage{subfloat}
\usepackage[draft]{fixme}
\usepackage{bm} 
\usepackage{bbm} 
\usepackage{epigraph}
\usepackage{endnotes}
\usepackage{ifpdf}
\usepackage{array}

\usepackage{multicol}
\usepackage{multirow}
\usepackage{graphicx}
\usepackage{tikz}

\newcommand{\A}{\mathbb{A}}
\newcommand{\Q}{\mathbb{Q}}
\newcommand{\N}{\mathbb{N}}
\newcommand{\Z}{\mathbb{Z}}
\newcommand{\R}{\mathbb{R}}
\newcommand{\F}{\mathbb{F}}
\newcommand{\C}{\mathbb{C}}
\renewcommand{\P}{\mathbb{P}}
\newcommand{\I}{\mathcal{I}}
\newcommand{\V}{\mathcal{V}}
\newcommand{\cH}{\mathcal{H}}
\newcommand{\ol}{\overline}
\newcommand{\J}{\mathcal{J}}
\newcommand{\X}{\mathcal{X}}
\theoremstyle{definition}
\newtheorem{prob}{Problem}
\newtheorem{extraProb}{Extra Problem}

\DeclareMathOperator{\conv}{conv}
\DeclareMathOperator{\cone}{cone}
\DeclareMathOperator{\ext}{ext}



% https://sites.math.washington.edu/~vinzant/teaching/380/
\begin{document}
 \thispagestyle{empty}
\begin{center} {\large \textbf{Math 380 -- Homework  4}} \smallskip \\ 
Due on Friday, May 8\ \\ \ \\
 \end{center}

You are welcome to talk with other students in the class about problems but should write up solutions on your own. 
Solutions can be handwritten or typed but need to be legible and submitted via Gradescope by the end of the day 
on Friday. You should justify all your answers in order to receive full credit.  
\bigskip 

\begin{flushleft}

{\bf Good practice problems} (not to be turned in): \\

CLO Section 2.5, Exercises 2, 10, 13,  16\\
CLO Section 2.6, Exercises 3, 5, 7, 9\\
CLO Section 2.7, Exercises 5, 10, 11\\
CLO Section 2.10, Exercises 12, 13\\ \ \\
\bigskip 




\begin{prob} CLO \S2.5, Exercises 14 and 15\\
\end{prob}
\bigskip 
\paragraph{\S2.5 Exercise 14} Let $f_1,f_2,\dots \in k[x_1,\dots,x_n]$ be an infinite sequence of polynomials. We can create an ascending chain of ideals:
\[\langle f_1\rangle \subseteq \langle f_1,f_2\rangle \subseteq \langle f_1,f_2,f_3\rangle \subseteq \dots \]
However, since polynomial rings are Noetherian, this chain must terminate. Denote $I_i = \langle f_1,\dots f_i\rangle$, and let $I = \bigcup_{i=1}^\infty I_i$. We have that $I$ is an ideal, since $f, g\in I$ implies that $f \in I_j$ and $g \in I_k$ for some $j,k \in \mathbb N$, and $f,g\in I_{\max\{j,k\}}\subseteq I$ (i.e. $I$ is the colimit). By Hilbert's Basis Theorem, $I$ must be finitely generated by some $\langle g_1,\dots, g_s \rangle$. Because each $g_i \in I$, we know that $g_i \in I_{(n_i)}$ for some $n_i \in \mathbb N$. Define $N = \max \{ n_i\}$, and we have that $g_i \in I_{N}$ for all $1 \le i \le s$. Thus, $I \subseteq I_N$, and by definition we also have $I_N \subseteq I$ and so $I_N = I$. Thus, for all $i \ge N+1$, we must have that $f_i \in I=I_N = \langle f_1,\dots, f_N\rangle$.

\vspace{1em}

\paragraph{\S2.5 Exercise 15} Let $f_1,f_2,\dots \in k[x_1,\dots,x_n]$ be an infinite sequence of polynomials, and consider $\mathbf V(f_1,f_2,\dots)$ as the solution set to $f_1=f_2=\dots = 0$. By the previous problem, we know that there is some $N \in \mathbb N$ such that $f_i \in \langle f_1,\dots f_N\rangle$ for all $i \in \mathbb N$. Denote this ideal $I = \langle f_1,\dots , f_N\rangle$, and note that $\mathbf V(I) \supseteq \mathbf V(f_1,f_2,\dots)$ as any point that is a solution to the infinite collection of equations is also a solution to a finite subset $I$. We can show that in fact $\mathbf V(I)\subseteq \mathbf V(f_1,f_2,\dots)$: let $a \in \mathbf V(I)$ be any point, and consider $f_i \in k[x_1,\dots,x_n]$ for any $i \in \mathbb N$. Since $f_i \in I$, we know that $f_i(a)=0$ and so $a$ is a solution for all $f_i$ and thus $a \in \mathbf V(f_1,f_2,\dots)$. Thus, $\mathbf V(f_1,\dots, f_N) = \mathbf V(f_1,f_2,\dots)$.

\newpage{}
\begin{prob} CLO \S2.6, Exercise 1 and 13\smallskip \\
\end{prob}
\smallskip

\paragraph{\S2.6 Exercise 1} Fix a monomial ordering, and let $I\subseteq k[x_1,\dots,x_n]$. Let $f \in k[x_1,\dots,x_n]$.
\begin{enumerate}
    \item[a.] By Hilbert's Basis theorem/Dickson's lemma, we can find a Gr\"obner basis for $I$, a finite set $G\subseteq I$ where $\text{LT}(I)=\langle \text{LT}(G)\rangle$. By the multivariable division algorithm, dividing $f$ by this Gr\"obner basis yields $f =r+ \sum q_i g_i$ for $g_i \in G$, where no term of $r$ is divisible by the leading term of any $g_i$. Let $g = \sum q_i g_i$, and since $g_i \in G \subseteq I$ note that $g \in I$. Let $c\cdot x^\alpha \in \text{LT}(I)=\langle \text{LT}(G)\rangle$ be any leading term appearing in $I$. Because $\langle \text{LT}(G)\rangle$ is a monomial ideal, $c\cdot x^\alpha$ must be divisible by some $\text{LT}(g_i)$ for $g_i \in G$. Assume FTSOC that $c\cdot x^\alpha$ divides some term of $r$. Then since some $\text{LT}(g_i)$ divides $c\cdot x^\alpha$, this implies $\text{LT}(g_i)$ divides that term of $r$, contradicting the property of multivariable division. Thus, no term $c\cdot x^\alpha \in \text{LT}(I)$ can divide any term of $r$, and we can write $f=g+r$, where $g \in I$ and no term of $r$ is divisible by any $\text{LT}(I)$. 
    \item[b.] Let $f = g+r=g'+r'$ with the properties above, where $g,g' \in I$ and no term of $r$ or $r'$ is in the ideal $\text{LT}(I)$. Then $f-f=g+r-g'-r'=0$ and $g-g'=r-r'$. Since $g,g'\in I$, we must have $r-r'\in I$. Since every monomial term of $r-r'$ must either have been a term of $r$ or $r'$, no nonzero term $x^\alpha$ appearing in $r-r'$ is divisible by any $\text{LT}(g_i)$. However, since $\text{LT}(r-r') \in \text{LT}(I)=\langle \text{LT}(G)\rangle$, some $\text{LT}(g_i)$ must divide $\text{LT}(r-r')$ as $\langle \text{LT}(G)\rangle$ is a monomial ideal. Thus, $\text{LT}(r-r')$ must be zero, and $r-r'=0$ since $0$ is the least element in every monomial order. Therefore, $r=r'$ and $g=g'$, showing uniqueness.
\end{enumerate}

\newpage{}
\paragraph{\S2.6 Exercise 13} Let $I \subseteq k[x_1,\dots,x_n]$ be an ideal with a Gr\"obner basis $G$. Consider $f,g\in k[x_1,\dots,x_n]$. By Exercise 1 above, we must have unique $f = q_f + \overline f^G$ and $g = q_g + \overline g^G$ where $q_f,q_g \in I$ and no term of $\overline f ^G$ or $\overline g^G$ is divisible by $\text{LT}(g_i)$ of any generator $g_i \in G$.
\begin{enumerate}
    \item[a.] 
    \[ \overline f^ G = \overline g^ G \iff f-g\in I.\]
    
    ($\implies$) Since $\overline f^G = \overline g^G$, we have that $f-g=q_f-q_g \in I$.

    ($\impliedby$) Since $f-g = (q_f - q_g) +\left(\overline f^G - \overline g^G \right)$ with $f-g \in I$ and $q_f - q_g \in I$, we must have $\overline f^G - \overline g^G \in I$. We can apply the same logic as part (b) of Exercise 1: every nonzero term of $\overline f^G-\overline g^G$ must not by divisible by any $\text{LT}(g_i)$, but $\text{LT}\left(\overline f^G -\overline g^G \right) \in \text{LT}(I)=\langle \text{LT}(G)\rangle$ so $\text{LT}\left(\overline f^G -\overline g^G \right)=0$ and in fact $\overline f^G - \overline g^G = 0$. Thus, $\overline f^G = \overline g^G$.
    \item[b.] %\[ \overline{f+g}^G = \overline{f}^G + \overline{g}^G. \]
    We can compute $f+g=(q_f + q_g)+ \left(\overline f^G+\overline g^G\right)$, and since every monomial term of $\overline f^G + \overline g^G$ must have come from $\overline f^G$ or $\overline g^G$, no term of $\overline f^G + \overline g^G$ is divisible by the leading term of any generator. Since $q_f+q_g \in I$ also, we know this is the unique decomposition with these properties, so $\overline{f+g}^G = \overline f^G + \overline g^G$ by definition.
    \item[c.] We can uniquely decompose $\overline f^G \overline g^G = q^\ast + \overline{\overline f^G \overline g^G}^G$ where $q^\ast \in I$ and no term of the remainder is divisible by the leading term of any generator. 
    Thus, we can write \begin{align*}
    fg &= \left(q_f + \overline f^G\right)\left(q_g + \overline g^G\right) = \left(q_f + \overline f^G\right)q_g + \overline g^G q_f +\overline f ^G \overline g^G \\ &= \left(q_f + \overline f^G\right)q_g + \overline g^G q_f +q^\ast +\overline{\overline f ^G \overline g^G}^G \end{align*}
    Since $q_f,q_g,q^\ast \in I$, we know that $\left(q_f + \overline f^G\right)q_g + \overline g^G q_f +q^\ast \in I$. Thus, we have found the unique decomposition of $fg$ into an element of $I$ plus a remainder with the desired properties, and $\overline{\overline f^G \overline g^G}^G = \overline{fg}^G$.
\end{enumerate}

\newpage{}
\begin{prob} CLO \S2.7, Exercises 14 \smallskip
 \\
Let $V = \{(a_1,b_1),\dots, (a_n,b_n)\} \subseteq k^2$ be a set of points where $a_1,\dots,a_n$ are all distinct, and define 
\[ h(x)=\sum_{i=1}^n b_i \prod_{j\ne i} \frac{x-a_j}{a_i-a_j} \in k[x]. \]
\begin{enumerate}
    \item[a.] For all $1 \le k \le n$, we have $h(a_k) = b_k$. We can evaluate
    \[h(a_k) = \sum_{i=1}^n b_i \prod_{j \ne i} \frac{a_k - a_j}{a_i - a_j} = \sum_{i=1}^n b_i \delta^i_k = b_k,\]
    where $\delta^i_k$ is the Kronecker delta. The product becomes the Kronecker delta because if $k \ne i$, then we have a term where $j =k$ and so $a_k-a_j = 0$. If $k=i$, then the product is taken over $\frac{a_i-a_j}{a_i-a_j} = 1$ ($j \ne i$ implies $a_i \ne a_j$ because $a_1,\dots,a_n$ are distinct). In the product, the maximum degree of $x$ that appears is $n-1$ as $j$ ranges over all the points except one. This is also the optimal upper bound, since we can pick $n$ points that force the function to cross $0$ exactly $n-1$ times, requiring a non-constant polynomial with $n-1$ roots.
    \item[b.] Let $g(x)$ be any polynomial with $\deg(g)\le n-1$ and $g(a_i)=b_i$ for all $1 \le i \le n$. Since $h$ has these same properties, we have that $(g-h)(a_i)=0$ for all $i$, and $\deg(g-h)\le \max\{\deg(g),\deg(h)\}\le n-1$. Thus, $g-h$ is a polynomial of degree at most $n-1$, but $g-h$ has $n$ roots. If $g-h$ were any non constant polynomial, $g-h$ would have at most $n-1$ roots, so $g-h$ is in fact the constant polynomial $0$. Since $g-h=0$, $g=h$, and we have shown that $h$ is the unique polynomial with these properties.
    \item[c.] Define $f = \prod_{i=1}^n (x-a_i) \in k[x]\subseteq k[x,y]$. We wish to show that $\langle f, y-h\rangle = \mathbf I(V)$. First, to show that $\langle f, y-h\rangle \subseteq \mathbf I(V)$, it suffices to show that $f,y-h\in \mathbf I(V)$. Let $p=(a_i,b_i)\in V$ for some $1 \le i \le n$. Then by construction, $f(p)=f(a_i)=0$, and similarly $(y-h)(p)=b_i-h(a_i)=0$.
    
    Now to show the reverse containment, let $g \in \mathbf I(V)$ be any polynomial. Fix the lexicographical order with $y > x$ and divide by $(f,y-h)$. By the multivariate division algorithm, we can write $g = q_1 \cdot f + q_2 \cdot (y-h) + r$ where no term of $r$ is divisible by any leading term of a divisor. Since $g,f,y-h \in \mathbf I(V)$, we know that $r \in \mathbf I(V)$ and thus $r(p) = 0$ for all $p \in V$. Because $|V| = n$, we know $r$ has $n$ distinct roots. 
    
    Under our monomial order, $\text{LT}(y-h) = y$ and $\text{LT}(f)=x^n$. Thus, every monomial term of $r$ is of the form $x^i$ where $i<n$, and so $r \in k[x]$ with $\deg(r) < n$. However, a non constant polynomial with $\deg r < n$ has at most $n-1$ distinct roots, so $r$ must in fact be the constant polynomial $0$. Thus, since $g = q_1 \cdot f + q_2 \cdot (y-h)$, $g \in \langle f, y-h\rangle$.
    \item[d.] With the $y>x$ lexicographical order, we can compute the $S$-polynomial $S(f,y-h)$. However, $S(f,y-h) \in \langle f,y-h\rangle=\mathbf I(V)$, and in part (c) above we have shown that for all polynomials in the ideal, the remainder upon division by $(f,y-h)$ is always $0$. Thus, $\overline{S(f,y-h)}^G = 0$ and by Buchberger's criterion $\{f,y-h\}$ is a Gr\"obner basis. We can also check that $\text{LC}(f) = \text{LC}(y-h)=1$, and neither $\text{LT}(f)=x^n$ nor $\text{LT}(y-h)=y$ divide the other, so we have no $\text{LT}(g_i) \in \langle \text{LT}(G \setminus \{g_i\})\rangle$. Thus, our set is a reduced Gr\"obner basis.
\end{enumerate}

\end{prob}
\bigskip 




\end{flushleft}
\end{document}